\subsection{PASSAGE FROM LAWS OF INFINITESIMAL OPERATION TO LAWS OF OPERATION}

PROPOSITION 11. Let \(\mathrm{G}_1\) and \(\mathrm{G}_2\) be Lie group germs and \(\mathbf{X}_1\) and \(\mathbf{X}_2\) manifolds of class \(\mathrm{C}^r\) ( \(r \geqslant 2\) ). For \(i = 1, 2\) , let \(\psi_i\) be a law chunk of left operation of class \(\mathrm{C}^r\) of \(\mathrm{G}_i\) on \(\mathbf{X}_i\) and \(\mathrm{D}_i\) the associated law of infinitesimal operation. Let \(\mu: \mathrm{G}_1 \to \mathrm{G}_2\) be a morphism and \(\phi: \mathrm{X}_1 \to \mathrm{X}_2\) a mapping of class \(\mathrm{C}^r\) . Suppose that, for all \(a \in \mathrm{L}(\mathrm{G})\) , the vector fields \((\mathrm{D}_{1})_a\) and \((\mathrm{D}_{2})_{\mathrm{L}(\mu)a}\) are \(\phi\) -related (Differentiable and Analytic Manifolds, R, 8.2.6). Then there exists a neighbourhood \(\Omega\) of \(\{e\} \times \mathbf{X}_1\) in \(\mathrm{G}_1 \times \mathbf{X}_1\) such that \(\phi(\psi_1(g,x)) = \psi_2(\mu(g),\phi(x))\) for all \((g,x) \in \Omega\) .

Let \(p_1: \mathrm{G}_1 \times \mathrm{X}_2 \to \mathrm{G}_1\) , \(p_2: \mathrm{G}_1 \times \mathrm{X}_2 \to \mathrm{X}_2\) be the canonical projections. For all \((g_1, x_2) \in \mathrm{G}_1 \times \mathrm{X}_2\) , let \(f_{g_1, x_2}\) be the mapping \(g_1 a \mapsto (\mathrm{D}_2)_{\mathrm{L}(\mu)a}(x_2)\) of \(\mathrm{T}_{g_1}(\mathrm{G}_1) = g_1\mathrm{L}(\mathrm{G}_1)\) into \(\mathrm{T}_{x_2}(\mathrm{X}_2)\) . The \(f_{g_1, x_2}\) define a morphism of \(p_1^* \mathrm{T}(\mathrm{G}_1)\) into \(p_2^* \mathrm{T}(\mathbf{X}_2)\) .

Let \(x_0 \in X_1\) . There exist an open neighbourhood G of e in \(G_1\) and an open neighbourhood X of \(x_0\) in \(X_1\) such that \(\psi_1(g, x)\) and \(\psi_2(\mu(g), \phi(x))\) are defined for \((g, x) \in G \times X\) . We write, for \((g, x) \in G \times X\) ,

\[
\alpha (g, x) = \phi (\psi_ {1} (g, x)) \in \mathrm{X} _ {2}, \quad \beta (g, x) = \psi_ {2} (\mu (g), \phi (x)) \in \mathrm{X} _ {2}.
\]

If \(\mathbf{G}\) and \(\mathbf{X}\) are sufficiently small, then, for all \((a,g,x)\in \mathrm{L}(\mathrm{G}_1)\times \mathrm{G}\times \mathrm{X},\)

\[
\begin{array}{r l} (\mathrm{T} \alpha) (a g, 0 _ {x}) & = (\mathrm{T} \phi) ((\mathrm{D} _ {1}) _ {a} (\psi_ {1} (g, x))) \\ & = (\mathrm{D} _ {2}) _ {L (\mu) a} (\phi (\psi_ {1} (g, x))) \\ & = (\mathrm{D} _ {2}) _ {L (\mu) a} \alpha (g, x), \end{array}
\]

\[
\begin{array}{r l} (\mathrm{T} \beta) (a g, 0 _ {x}) = & (\mathrm{T} \psi_ {2}) (\mathrm{L} (\mu) a. \mu (g), \phi (x)) \\ = & (\mathrm{D} _ {2}) _ {\mathrm{L} (\mu) a} (\psi_ {2} (\mu (g), \phi (x))) \\ = & (\mathrm{D} _ {2}) _ {\mathrm{L} (\mu) a} \beta (g, x). \end{array}
\]

Hence for \(x \in \mathbf{X}\) , the morphisms \(g \mapsto \alpha(g, x)\) and \(g \mapsto \beta(g, x)\) are integrals of \(f\) ; as

\[
\beta (e, x) = \phi (x) = \alpha (e, x)
\]

for all \(x \in \mathbf{X}\) , it follows from Differentiable and Analytic Manifolds, R, 9.3.7, that \(\alpha\) and \(\beta\) coincide on a neighbourhood of \((e, x_0)\) . Hence the proposition.

COROLLARY. Let G be a Lie group germ and X a manifold of class \(\mathbf{C}^{r}\) . Consider two law chunks of left operation of class \(\mathbf{C}^{r}\) of G on X. Suppose that, for all \(a \in L(G)\) , the corresponding vector field \(D_{a}\) on X is the same for both law chunks. Then these two law chunks coincide on a neighbourhood of \(\{e\} \times X\) .

THEOREM 6. Let G be a Lie group germ, X a manifold of class \(\mathbf{C}^r\) ( \(r \geqslant 2\) ) and \(x_0\) a point of X. Let \(a \mapsto \mathrm{D}_a\) be a law of left infinitesimal operation of class \(\mathbf{C}^{r-1}\) of \(\mathrm{L(G)}\) on X.

\begin{enumerate}
\def\labelenumi{(\roman{enumi})}
\item
  There exists an open neighbourhood \(\mathbf{X}'\) of \(x_0\) in \(\mathbf{X}\) and a law chunk of left operation of class \(\mathbf{C}^{r-1}\) of \(\mathbf{G}\) on \(\mathbf{X}'\) such that the associated law of infinitesimal operation is \(a \mapsto \mathbf{D}_a|\mathbf{X}'\) .
\item
  Let there be two law chunks of left operation of class \(\mathbf{C}^{r-1}\) of \(\mathbf{G}\) on an open neighbourhood \(\mathbf{X}''\) of \(x_0\) ; if they admit \(a \mapsto \mathbf{D}_a|\mathbf{X}''\) as associated law of infinitesimal operation, they coincide on a neighbourhood of \((e, x_0)\) .
\end{enumerate}

Assertion (ii) follows from the Corollary to Proposition 11. We prove (i). For all \((g, x) \in \mathbf{G} \times \mathbf{X}\) and all \(a \in \mathrm{L}(\mathbf{G})\) , we write

\[
\mathrm{Q} _ {a} (g, x) = (a g, \mathrm{D} _ {a} (x)) \in \mathrm{T} _ {g} (\mathrm{G}) \times \mathrm{T} _ {x} (\mathrm{X}).
\]

Let \(S_{(g,x)}\) be the set of \(Q_{a}(g,x)\) for \(a\in L(G)\) . By Lemma 5 of no. 6, the \(S_{(g,x)}\) are the fibres of a vector subbundle S of \(T(G)\times T(X)\) . Let a, b be in \(L(G)\) ; then

\[
\begin{array}{r l} {[ \mathrm{Q} _ {a}, \mathrm{Q} _ {b} ] (g, x) = ([ \mathrm{R} _ {a}, \mathrm{R} _ {b} ] (g), [ \mathrm{D} _ {a}, \mathrm{D} _ {b} ] (x))} \\ & = (- \mathrm{R} _ {[ a, b ]} (g), - \mathrm{D} _ {[ a, b ]} (x)) \\ & = \mathrm{Q} _ {- [ a, b ]} (g, x) \end{array}\tag{§ 3, 18.6}
\]

and hence S is integrable (Differentiable and Analytic Manifolds, R, 9.3.3, (iv)). By Differentiable and Analytic Manifolds, R, 9.3.7, there exist an open neighbourhood \(\mathrm{G}_1\) of \(e\) in G, an open neighbourhood \(\mathbf{X}_1\) of \(x_0\) in X and a mapping \((g,x)\mapsto gx\) of class \(\mathrm{C}^{r - 1}\) of \(\mathrm{G}_1\times \mathrm{X}_1\) into X such that \(ex = x\) for all \(x\in \mathbf{X}_1\) and

\[
(a g) x = \mathrm{D} _ {a} (g x) \quad \text { for } a \in \mathrm{L} (\mathrm{G}), g \in \mathrm{G} _ {1}, x \in \mathrm{X} _ {1}.\tag{6}
\]

In particular

\[
a x = \mathrm{D} _ {a} (x).\tag{7}
\]

Let \(G_{2}\) be an open neighbourhood of e in \(G_{1}\) and \(X_{2}\) an open neighbourhood of \(x_{0}\) in \(X_{1}\) such that \(gg'\) is defined and belongs to \(G_{1}\) for g, \(g'\) in \(G_{2}\) and gx is defined and belongs to \(\mathbf{X}_1\) for \((g,x)\in \mathbf{G}_2\times \mathbf{X}_2\) . Consider the mappings \(\alpha_{1},\alpha_{2}\) of \(\mathrm{G}_2\times (\mathrm{G}_2\times \mathrm{X}_2)\) into X defined by

\[
\alpha_ {1} (g, (h, x)) = g (h x), \quad \alpha_ {2} (g, (h, x)) = (g h) x.
\]

They are of class \(\mathbf{C}^{r - 1}\) . Then

\[
\alpha_ {1} (e, (h, x)) = h x = \alpha_ {2} (e, (h, x)).
\]

On the other hand

\[
\begin{array}{r l} \mathrm{T} (\alpha_ {1}) (a g, 0 _ {(h, x)}) & = (a g) (h x) \\ & = \mathrm{D} _ {a} (g (h x)) \qquad \text {by (6)} \\ & = \mathrm{D} _ {a} (\alpha_ {1} (g, (h, x))), \\ \mathrm{T} (\alpha_ {2}) (a g, 0 _ {(h, x)}) & = (a g h) x \\ & = \mathrm{D} _ {a} ((g h) x) \qquad \text {by (6)} \\ & = \mathrm{D} _ {a} (\alpha_ {2} (g, (h, x))). \end{array}
\]

By Differentiable and Analytic Manifolds, R, 9.3.7, \(\alpha_{1}\) and \(\alpha_{2}\) coincide on a neighbourhood of \((e, (e, x_{0}))\) . Then (i) follows from (7) and Proposition 23 of § 1, no. 11.

COROLLARY 1. Let G be a Lie group germ and X a paracompact manifold of class \(\mathbf{C}^r\) ( \(r \geqslant 2\) ). Let \(a \mapsto \mathrm{D}_a\) be a law of left infinitesimal operation of class \(\mathbf{C}^{r-1}\) of L(G) on X.

\begin{enumerate}
\def\labelenumi{(\roman{enumi})}
\item
  There exists a law chunk of left operation of class \(\mathbf{C}^{r-1}\) of \(\mathbf{G}\) on \(\mathbf{X}\) such that the associated law of infinitesimal operation is \(a \mapsto \mathbf{D}_a\) .
\item
  Two laws of left operation of class \(\mathbf{C}^{r-1}\) of G on X which admit \(a \mapsto \mathbf{D}_a\) as associated law of infinitesimal operation coincide on a neighbourhood of \(\{e\} \times \mathbf{X}\) .
\end{enumerate}

Assertion (ii) follows from the Corollary to Proposition 11. By Theorem 6 (i), there exist an open covering \((\mathbf{X}_i)_{i\in I}\) of \(\mathbf{X}\) and, for all \(i\in I\) , a law chunk of left operation \(\psi_i\) of class \(\mathbf{C}^{r-1}\) of \(\mathbf{G}\) on \(\mathbf{X}_i\) such that the associated law of infinitesimal operation is \(a\mapsto \mathrm{D}_a|\mathbf{X}_i\) . As \(\mathbf{X}\) is paracompact, the covering \((\mathbf{X}_i)_{i\in I}\) can be assumed to be locally finite. For all \((i,j)\in I\times I\) and all \(x\in \mathbf{X}_i\cap \mathbf{X}_j\) , \(\psi_i\) and \(\psi_j\) coincide on a neighbourhood of \((e,x)\) (Corollary to Proposition 11). As \(\mathbf{X}\) is normal, we can apply Proposition 24 of § 1, no. 11, which proves (i).

COROLLARY 2. Let X be a paracompact manifold of class \(\mathbf{C}^r\) ( \(r \geqslant 2\) ) and \(\xi\) a vector field of class \(\mathbf{C}^{r-1}\) on X. There exists a law chunk of operation \(\psi\) of class \(\mathbf{C}^{r-1}\) of K on X such that for all \(x \in \mathbf{X}\) , \(\xi(x)\) is the image under \(t \mapsto (t, x)\) of the tangent vector 1 to K at 0. Two law chunks of operation with the above property coincide on a neighbourhood of \(\{0\} \times \mathbf{X}\) .

This is a special case of Corollary 1.

Remark. Laws of left operation can of course be replaced, throughout this no., by laws of right operation.

